%%%%%%%%%%%%%%%%%%%%%%%%%%%%%%%%%%%%%%%%%%%%%%%%%%%%%%%%%%%%%%%%%%%%%%%%
% Plantilla TFG/TFM
% Universidad de Murcia. Facultad de Informática
% Realizado por: José Manuel Requena Plens
% Modificado: Pablo José Rocamora Zamora
% Contacto: pablojoserocamora@gmail.com
%%%%%%%%%%%%%%%%%%%%%%%%%%%%%%%%%%%%%%%%%%%%%%%%%%%%%%%%%%%%%%%%%%%%%%%%


\chapter*{Resumen}

La proliferación de desinformación en el ecosistema digital ha convertido la verificación automatizada de hechos en una necesidad urgente. Los Grandes Modelos de Lenguaje (LLMs) presentan capacidades notables para el procesamiento del lenguaje natural, pero su tendencia a generar información plausible pero falsa (alucinación) los hace poco fiables para operar de forma aislada en tareas donde la veracidad es crítica. A esto se suma la degradación temporal de su conocimiento interno, que permanece estático desde el momento en que concluye el entrenamiento.

Este trabajo propone el diseño, implementación y validación de un sistema de verificación de noticias basado en una arquitectura de Generación Aumentada por Recuperación (RAG). El sistema integra un motor de búsqueda Elasticsearch sobre un corpus de 490.956 noticias, construido tras un proceso de extracción y limpieza de datos a partir de 804.083 registros brutos. Se implementan y comparan dos estrategias de recuperación: búsqueda léxica mediante el algoritmo BM25 y búsqueda semántica vectorial mediante el modelo de \textit{embeddings} BGE-M3. El componente generador fue evaluado con tres modelos de distinta escala: Phi-3 (SLM local), Llama 3.3 70B y Qwen 3 32B, los dos últimos accedidos vía API.

La validación se realizó sobre un \textit{dataset} de 220 preguntas clasificadas en factuales, de inferencia y trampa, usando el paradigma \textit{LLM-as-a-Judge} complementado con supervisión humana. Los resultados demuestran que el sistema RAG óptimo —combinando recuperación semántica con Llama 3.3 70B— alcanza un 75,9\% de precisión y reduce la tasa de alucinación al 2,2\% de los errores cometidos. El trabajo concluye que la calidad del motor de recuperación es el factor más determinante del rendimiento final, y que el diseño de arquitecturas RAG representa actualmente el enfoque más seguro para aplicar inteligencia artificial generativa en entornos de verificación factual.

